% Language Lab tutorial written in LaTeX. The hub (index.html) reads this subset directly:
%   \tutorialid \title \group \summary \roles \task  (metadata)
%   \begin{flow}{lanes} \node{id}{lane}{type}{text}{next[:Label],...}{optional link} \end{flow}
%   \section[role]{Title}\label{id}  then plain text, enumerate (steps), itemize (list),
%   \begin{script}...\end{script} (what to say), \begin{note}...\end{note}
% Roles: student, la, lc, dept, all.  Node types: start, step, decision, end, email.
% Inline: \textbf \emph \texttt \href{url}{text}; {email} becomes the chosen campus's Lab email.
% To use it: save as t/<id>.tex and put "t/<id>.tex" in tutorials.json's "tutorials" list.
\documentclass{article}
\newcommand{\tutorialid}[1]{}\newcommand{\group}[1]{}\newcommand{\summary}[1]{\emph{#1}}
\newcommand{\roles}[1]{}\newcommand{\task}[2]{}\newenvironment{flow}[1]{}{}\newcommand{\node}[6]{}
\newenvironment{script}{\begin{quote}}{\end{quote}}\newenvironment{note}{\begin{quote}\itshape}{\end{quote}}
\begin{document}
\tutorialid{example}
\title{Example tutorial}
\group{desk}
\summary{One line that says what this tutorial is for.}
\roles{la,lc}
\task{la}{What the lab assistant does}
\task{lc}{What the coordinator does}

\begin{flow}{student,la,lc}
\node{a}{student}{start}{Student asks}{b}
\node{b}{la}{decision}{Can you settle it now?}{c:Yes,d:No}
\node{c}{la}{end}{Settle it}{}{\#do}
\node{d}{student}{email}{Emails {email}}{e}
\node{e}{lc}{end}{Coordinator answers}{}
\end{flow}

\section[la]{Do it}\label{do}
\begin{enumerate}
  \item Open \texttt{the-site.ung.edu}.
  \item Click \textbf{Save}.
\end{enumerate}
\begin{note}Not sure? Send it to email.\end{note}

\section[lc]{Coordinator}\label{lc}
\begin{itemize}
  \item Answer by email.
\end{itemize}
\end{document}
